\documentclass[a4paper,12pt]{jarticle}

\usepackage[dvipdfmx]{geometry}
\usepackage{longtable}
\usepackage{array}
\usepackage{enumitem}

\geometry{margin=25mm}

\title{要求仕様書\\献立提案システム}
\author{}
\date{}

\begin{document}

\maketitle

\section{システム概要}

本システムは、日々の献立を考える負担を軽減することを目的としたWebシステムである。
ユーザが現在所有している食材を入力すると、その食材を参考にしたおすすめの献立を提案する。

なお、本システムは段階的な機能拡張を前提として開発を行う。

\section{機能要求}

\subsection{FR-01 食材入力機能}

\subsubsection*{概要}
ユーザは使用可能な食材を自由入力できる。

\subsubsection*{入力}
\begin{itemize}
    \item 食材名（複数入力可）
    \item 入力形式は自由入力
\end{itemize}

\subsubsection*{出力}
入力内容を献立検索機能へ受け渡す。

\subsubsection*{将来拡張}
\begin{itemize}
    \item 入力補完
    \item 食材候補表示
    \item 音声入力
    \item 画像認識による入力
\end{itemize}

\subsection{FR-02 献立検索機能}

\subsubsection*{概要}
入力された食材をもとにおすすめ献立を検索する。

\subsubsection*{条件}
\begin{itemize}
    \item 入力食材の一部が使用されていれば候補とする。
    \item 入力食材をすべて使用する必要はない。
\end{itemize}

\subsubsection*{将来拡張}
\begin{itemize}
    \item 使用率順で検索
    \item 調理時間による絞り込み
    \item カロリーによる絞り込み
    \item 栄養バランスによる検索
\end{itemize}

\subsection{FR-03 献立提案機能}

\subsubsection*{概要}
検索結果からおすすめ献立を表示する。

\subsubsection*{表示件数}
3件

\subsubsection*{表示内容}
\begin{itemize}
    \item 献立名
\end{itemize}

\subsubsection*{将来拡張}
\begin{itemize}
    \item レシピ表示
    \item 材料表示
    \item 調理時間表示
    \item 栄養情報表示
    \item 人気順表示
\end{itemize}

\subsection{FR-04 献立種類選択機能}

\subsubsection*{概要}
検索前に表示する献立の種類を選択できる。

\subsubsection*{選択項目}
\begin{itemize}
    \item 献立一式
    \item 主菜
    \item 副菜
    \item 汁物
\end{itemize}

\subsubsection*{将来拡張}
\begin{itemize}
    \item デザート
    \item お弁当
    \item 朝食
    \item 作り置き料理
\end{itemize}

\subsection{FR-05 エラー表示機能}

\subsubsection*{概要}
入力された食材名が認識できない場合、エラーメッセージを表示する。

\subsubsection*{エラー例}
\begin{itemize}
    \item 登録されていない食材です。
    \item 食材を入力してください。
\end{itemize}

\subsubsection*{将来拡張}
\begin{itemize}
    \item 類似食材候補の提示
    \item 誤字修正候補
    \item 入力補完
\end{itemize}

\section{非機能要求}

\subsection{NFR-01 利用環境}

Webブラウザ上で利用できること。

\subsection{NFR-02 操作性}

誰でも直感的に利用できる画面構成とする。

\subsection{NFR-03 応答性能}

検索開始から検索結果表示までを短時間で完了できること。

※具体的な応答時間は今後決定する。

\subsection{NFR-04 保守性}

機能追加・仕様変更が容易な構成とする。

\subsection{NFR-05 拡張性}

以下のような機能追加が可能な設計とする。

\begin{itemize}
    \item レシピ表示
    \item AIによる献立提案
    \item 栄養管理
    \item アレルギー対応
    \item 人数指定
    \item 冷蔵庫管理
    \item お気に入り登録
    \item 履歴表示
\end{itemize}

\section{入力仕様}

\begin{longtable}{|p{4cm}|p{9cm}|}
\hline
項目 & 内容\\
\hline
入力方式 & 自由入力\\
\hline
入力数 & 複数\\
\hline
必須入力 & 食材名\\
\hline
入力例 & 鶏肉、玉ねぎ、卵\\
\hline
\end{longtable}

\section{出力仕様}

\begin{longtable}{|p{4cm}|p{9cm}|}
\hline
項目 & 内容\\
\hline
提案件数 & 3件\\
\hline
表示内容 & 献立名\\
\hline
\end{longtable}

\section{エラー仕様}

\begin{longtable}{|p{6cm}|p{7cm}|}
\hline
条件 & 処理\\
\hline
食材未入力 & エラーメッセージ表示\\
\hline
存在しない食材 & エラーメッセージ表示\\
\hline
\end{longtable}

\section{システム構成}

\begin{center}
利用者

$\downarrow$

食材入力

$\downarrow$

献立種類選択

$\downarrow$

検索

$\downarrow$

おすすめ献立表示
\end{center}

\section{今後の拡張方針}

本システムは段階的な開発を前提とする。

\subsection*{初期バージョン}

\begin{itemize}
    \item 食材入力
    \item 献立種類選択
    \item 献立検索
    \item 献立名表示
    \item エラー表示
\end{itemize}

\subsection*{将来的な追加機能}

\begin{itemize}
    \item 入力候補表示
    \item レシピ表示
    \item 栄養情報表示
    \item AIによる献立提案
    \item アレルギー対応
    \item 人数指定
    \item 冷蔵庫管理
    \item 履歴機能
    \item お気に入り機能
\end{itemize}

\end{document}